\chapter{The Gaussian Binary Perceptron (Talagrand, Chapter~2)}\label{chap:sg-talagrand-perceptron}

This chapter starts Chapter~2 of Michel Talagrand's \emph{Spin Glasses: A Challenge for Mathematicians}
(Volume~I, Springer, 2010) with Section~2.1, the Gaussian binary perceptron. It is kept apart from
Chapter~\ref{chap:sg-talagrand}, which treats Sections~1.1--1.14 of the same volume, but uses the same
conventions: spins are functions $\mathrm{Fin}\,N\to\mathbb Z^\times$ and the disorder is a product of
standard Gaussian laws. Each statement below begins with a bold tag \textbf{[Talagrand, \dots]} naming the
passage of the book it corresponds to.

\section{The Gaussian binary perceptron: fixed-spin fields and the first moment (Talagrand Section~2.1)}\label{sec:sg-tal:perceptron-introduction}

Section~2.1 of Talagrand's Volume~I introduces the binary perceptron. There are $M$ random linear
constraints on $N$ spins; a spin configuration is \emph{feasible} when it satisfies all of them. This
section formalizes the exact fixed-spin laws behind the book's remark that the constraint fields are
``of order one'', and the first-moment argument showing that no feasible configuration exists when
$M/N$ tends to a limit above $1$.

\medskip\noindent\textbf{Notation.} Write $\Sigma_N=\{-1,+1\}^N$ and let $g=(g_{k,i})_{k<M,\,i<N}$ be
independent standard Gaussians. The book writes $g_{i,k}$; the Lean statements put the index $k$ of the
independent constraint first, so the column $k$ is the vector $(g_{k,i})_{i<N}$. For $\sigma\in\Sigma_N$ put
\[
S_k(g,\sigma)=\frac1{\sqrt N}\sum_{i<N}g_{k,i}\sigma_i,\qquad
A(g)=\Bigl\{\sigma\in\Sigma_N:\ \sum_{i<N}g_{k,i}\sigma_i\ge0\ \text{for all }k<M\Bigr\},\qquad
C(g)=|A(g)|.
\]
The closed half-space $\ge0$ is used, as in Eq.~(2.1). The sign of $\sum_i g_{k,i}\sigma_i$ and that of
$S_k(g,\sigma)$ agree because $\sqrt N>0$.

\begin{remark}[Scope of Section~2.1 in this blueprint]\label{rem:sg-tal:perceptron-scope}
\talsource{Section~2.1}
Formalized here: the exact laws of the constraint fields of \emph{fixed} spins (Lemma~\ref{lem:sg-tal:perceptron-fields}),
the expected number of feasible configurations, Eq.~(2.2), and the emptiness statement for $M/N\to\alpha>1$
(Theorems~\ref{thm:sg-tal:perceptron-first-moment} and~\ref{thm:sg-tal:perceptron-empty}).
Not formalized here: the comparison of Gaussian and spherical normal vectors in the opening paragraph; the
exponential-penalty limit (2.4)--(2.5); the energy (2.6) with smooth potential (2.7) and the pressure (2.11);
Research Problem~2.1.1, whose sharp threshold and entropy function are unknown in the book; the Volume~II
bound with $\alpha_0<1$; and the cavity equations (2.8), (2.10) of Sections~2.2--2.4. The first-moment bound
does not locate the threshold; it proves emptiness only above $\alpha=1$.
\end{remark}

\begin{lemma}[Constraint fields of fixed spins]\label{lem:sg-tal:perceptron-fields}
  \lean{ProbabilityTheory.map_dotProduct_pi_pi_gaussianReal_of_dotProduct_self,
        ProbabilityTheory.iIndepFun_dotProduct_pi_pi_gaussianReal,
        ProbabilityTheory.integral_normalized_gaussian_linear_form_mul,
        SpinGlasses.map_normalized_constraint_field_eq_gaussianReal,
        SpinGlasses.integral_normalized_constraint_field_eq_zero,
        SpinGlasses.integral_normalized_constraint_field_mul_eq,
        SpinGlasses.integrable_normalized_constraint_field_mul,
        SpinGlasses.hasGaussianLaw_normalized_constraint_fields,
        SpinGlasses.map_toLp_normalized_constraint_fields_eq_multivariateGaussian,
        SpinGlasses.iIndepFun_normalized_constraint_fields,
        SpinGlasses.map_normalized_constraint_fields_eq_pi_gaussianReal}
  \leanok
\talsource{Section~2.1, the remark that the fields are of order one}
Let $N\ge1$ and let $x=(x_i)_{i<N}$ be i.i.d. standard Gaussians. For fixed (deterministic) $\sigma,\tau\in\Sigma_N$
let $S(\sigma)=N^{-1/2}\sum_i x_i\sigma_i$.
\begin{enumerate}
\item $S(\sigma)$ has the law $\mathcal N(0,1)$.
\item $\mathbb E\,S(\sigma)=0$, and $S(\sigma)S(\tau)$ is integrable with
$\mathbb E\,[S(\sigma)S(\tau)]=N^{-1}\sum_{i<N}\sigma_i\tau_i=R(\sigma,\tau)$; the second statement holds for every $N$.
\item For finitely many fixed spins $(\sigma^a)_{a\in A}$ the vector $(S(\sigma^a))_a$ is jointly Gaussian,
centered, with covariance matrix $\bigl(R(\sigma^a,\sigma^b)\bigr)_{a,b}$. This matrix may be singular, for
example when two of the spins coincide.
\item For the disorder $g$ and one fixed spin $\sigma^k$ per column $k<M$, the fields
$N^{-1/2}\sum_i g_{k,i}\sigma^k_i$ are independent, and (for $N\ge1$) they are i.i.d. $\mathcal N(0,1)$.
\end{enumerate}
\end{lemma}
\begin{proof}
  \leanok
\emph{Step 1 (coefficients).} Put $a_i=\sigma_i/\sqrt N$. Since $\sigma_i^2=1$,
\[
\sum_{i<N}a_ib_i=\frac1N\sum_{i<N}\sigma_i\tau_i,\qquad \sum_{i<N}a_i^2=\frac1N\cdot N=1\quad(N\ge1),
\]
where $b_i=\tau_i/\sqrt N$. The second display is the self-overlap equal to $1$ on the hypercube.

\emph{Step 2 (Gaussian linear forms).} For i.i.d. standard Gaussians $x$ the linear form $a\cdot x$ is centered
with $\mathbb E[(a\cdot x)(b\cdot x)]=a\cdot b$, and $a\cdot x\sim\mathcal N(0,a\cdot a)$. With Step~1 this gives
(1) and (2). The family $(a^r\cdot x)_r$ over finitely many coefficient vectors is jointly Gaussian with Gram
matrix $(a^r\cdot a^s)$, which gives (3).

\emph{Step 3 (columns).} The columns $g_{k,\cdot}$ are independent, so functions of different columns are
independent. If each coefficient vector has square sum one, each field is $\mathcal N(0,1)$, and the law of the
vector of fields is the product law $\mathcal N(0,1)^{\otimes M}$. This gives (4).

The statements concern fixed spins under the disorder law; they do not say that the field stays standard normal
when $\sigma$ is drawn from a Gibbs measure that depends on $g$.
Formalized in \texttt{SpinGlasses/Talagrand/Perceptron/GaussianFields.lean}, with the general Gaussian
facts in \texttt{ForMathlib/Probability/Moments/GaussianGram.lean} and
\texttt{GaussianPiLinearForm.lean}.
\end{proof}

\begin{theorem}[Expected number of feasible configurations]\label{thm:sg-tal:perceptron-first-moment}
  \lean{ProbabilityTheory.gaussianReal_Ici_zero,
        SpinGlasses.measurableSet_feasible_configuration,
        SpinGlasses.card_feasible_configurations_eq_sum_indicator,
        SpinGlasses.measurable_card_feasible_configurations,
        SpinGlasses.measureReal_feasible_configuration_eq,
        SpinGlasses.integral_card_feasible_configurations_eq,
        SpinGlasses.measureReal_exists_feasible_configuration_le}
  \leanok
\talsource{Eq.~(2.2) and the paragraph after it}
For $N\ge1$ and every $M\ge0$:
\begin{enumerate}
\item for each fixed $\sigma\in\Sigma_N$, $\ \mathbb P(\sigma\in A(g))=(1/2)^M$;
\item $\mathbb E\,C(g)=2^N/2^M$, where the powers are real powers of natural numbers;
\item $\mathbb P(A(g)\ne\varnothing)\le 2^N/2^M$.
\end{enumerate}
The count $C$ and the events are measurable.
\end{theorem}
\begin{proof}
  \leanok
\emph{Step 1 (one constraint).} A standard Gaussian charges $[0,\infty)$ with mass $1/2$, since it is
symmetric and a point has no mass:
\[
\gamma([0,\infty))=\gamma((0,\infty))=\gamma((-\infty,0)),\qquad 2\,\gamma([0,\infty))=1.
\]

\emph{Step 2 (all constraints).} By Lemma~\ref{lem:sg-tal:perceptron-fields}(4) the fields of a fixed spin are
i.i.d. standard normal, so
\[
\mathbb P(\sigma\in A(g))=\gamma^{\otimes M}\bigl([0,\infty)^M\bigr)=(1/2)^M.
\]

\emph{Step 3 (expectation).} Write $C(g)=\sum_{\sigma}\mathbf 1_{\{\sigma\in A(g)\}}$. Each indicator is
bounded and measurable, so
\[
\mathbb E\,C=\sum_{\sigma\in\Sigma_N}\mathbb P(\sigma\in A)=2^N\cdot 2^{-M}=\frac{2^N}{2^M}.
\]
The quotient is written with real powers; the natural-number difference $N-M$ would be wrong for $M>N$.

\emph{Step 4 (Markov).} $C$ is nonnegative and integrable, and $C\ge1$ on $\{A\ne\varnothing\}$, hence
$\mathbb P(A\ne\varnothing)\le\mathbb E\,C$.

The hypothesis $N\ge1$ is needed: for $N=0$ every field is zero and the single empty configuration satisfies all
constraints, so $\mathbb P(\sigma\in A)=1\ne(1/2)^M$ when $M>0$. The case $M=0$ is allowed.
Formalized in \texttt{SpinGlasses/Talagrand/Perceptron/FirstMoment.lean}, with the half-line fact in
\texttt{ForMathlib/Probability/Distributions/Gaussian/Symmetry.lean}.
\end{proof}

\begin{theorem}[No feasible configuration above the first-moment threshold]\label{thm:sg-tal:perceptron-empty}
  \lean{SpinGlasses.tendsto_measureReal_exists_feasible_configuration_zero}
  \leanok
\talsource{The paragraph after Eq.~(2.2)}
Let $M:\mathbb N\to\mathbb N$ and $\alpha>1$ with $M(N)/N\to\alpha$. Then
\[
\mathbb P_{N,M(N)}\bigl(A(g)\ne\varnothing\bigr)\longrightarrow0\qquad(N\to\infty),
\]
where each probability is taken under the product Gaussian law of its own disorder $g\in\mathbb R^{M(N)\times N}$.
\end{theorem}
\begin{proof}
  \leanok
Put $\theta=(\alpha-1)/2>0$. Since $M(N)/N\to\alpha>1+\theta$, eventually $M(N)\ge(1+\theta)N$ and $N\ge1$.
By Theorem~\ref{thm:sg-tal:perceptron-first-moment}(3), for such $N$
\[
\mathbb P_{N,M(N)}(A\ne\varnothing)\le\frac{2^N}{2^{M(N)}}\le\frac{2^N}{2^{(1+\theta)N}}=\bigl(2^{-\theta}\bigr)^N.
\]
Since $2^{-\theta}<1$, the right-hand side tends to $0$, and probabilities are nonnegative. No coupling of
different $N$ is needed.
Formalized in \texttt{SpinGlasses/Talagrand/Perceptron/FirstMoment.lean}.
\end{proof}

\begin{remark}[Comparison with the book, Section~2.1]\label{rem:sg-tal:perceptron-comparison}
\talsource{Section~2.1}
\begin{itemize}
\item The book says the fields are ``of order one''. Lemma~\ref{lem:sg-tal:perceptron-fields} makes this exact:
a fixed spin gives a standard normal field. This is a statement for deterministic spins under the disorder
law; it fills in what the book leaves implicit, and does not apply to Gibbs-distributed spins.
\item The book computes $\mathbb E\,C=2^N/2^M$ and states emptiness for $\alpha>1$ in one sentence. The Lean
proof adds the explicit bound $\mathbb P(A\ne\varnothing)\le(2^{-\theta})^N$, uses real powers so that $M>N$ is
handled, and records that $N\ge1$ is used.
\item The index order $g_{i,k}$ of the book is reversed in Lean to $g_{k,i}$, so that the outer index
groups the independent constraints.
\item This is the book's first-moment argument. It neither computes the true threshold nor shows
nonemptiness below it, so it does not address Research Problem~2.1.1.
\end{itemize}
\end{remark}
